\documentclass[10pt]{article}

\usepackage[utf8]{inputenc}

\usepackage[T1]{fontenc}

\usepackage{amsmath}

\usepackage{amsfonts}

\usepackage{amssymb}

\usepackage[version=4]{mhchem}

\usepackage{stmaryrd}

\usepackage{bbold}

\usepackage{graphicx}

\usepackage[export]{adjustbox}

\graphicspath{ {./images/} }

\usepackage{mathrsfs}



\begin{document}

и пусть



\[

g^{\prime}= \begin{cases}g S_{k}, & n_{k}<0 \\ g S_{k}^{-1}, & n_{k}>0\end{cases}

\]



Тогда требуется, чтобы әлемент $g$ также принадлежал этой системе (әлемент $g^{\prime}$ можно представлять себе как редуцированное слово, полученное из $\mathrm{g}$ зачеркиванием его последней буквы). Элемент 1 принадлежит каждой IIl. $\mathbf{c}$.



Введена О. IIl рейером (O. Schreier) в 20-х гг., см. [1].



Лит.: [1] М а с с и у., С т о л л и н г с Д ж., Алгебраическая топология. Введение, пер. с англ., М., 1977. полво е многообрази е, - паракомпактное комплексное аналитическое многообразие $M$, обладающее следующими свойствами:



\begin{enumerate}

  \item для любого компакта $K \subset M$ множество

\end{enumerate}



\[

\left\{x \in X \| f(x)\left|\leqslant \sup _{z \in K}\right| f(z) \mid(f \in G(M))\right\},

\]



где $\sigma(M)$ - алгебра голоморфных функций на $M$, компактно (го ло мо рфвая в ы пу к лость);



\begin{enumerate}

  \setcounter{enumi}{1}

  \item для любых двух различных точек $x, y \in M$ существует такая функция $f \in G(M)$, что $f(x) \neq f(y)$ (го ломорфвая отделимость);



  \item в окрестности любой точки существует голоморфная карта, координатные функции к-рой принадлежат $\sigma(M)$.



\end{enumerate}



Условие голоморфной выпуклости можно заменить следующим: для любой последовательности точек $\left\{x_{n} \mid n=1\right.$, $2, \ldots\} \subset M$, не имеюшей предельных точек, существует такая фуунция $f \in G(M)$, что $\sup \left|f\left(x_{n}\right)\right|=\infty$.



Класс III. м. был введен в рассмотрение К. Штейном [1], как естественное обобщение голоморфиости областей в $\mathbb{C}^{n}$. Всякое замкнутое аналитич. Іодмногообразие в $\mathbb{C}^{n}$ является III. м.; обратно, любое $n$-мерное III. м. допускает собственное голоморфиое вложевие в $\mathbb{C}^{2 n}$. Всякая некомпактная риманова поверхность является Ш. м. Непосредственным обобщением Ш. м. являются III тейна пространства.



Jum.: [1] S t e i n K., "Math. Ann.», 1951, Bd 123, S. 201-22. См. также лит. при статье Штейна пространство.



IIITEЙНА ПРОСТР



IIITEЙНА ПРОСТРАНСТВО, г о л о м о $\mathrm{p}$ Ф н о п $о$ лн о е п р о с т р а н с т в о,- паракомпактное комплексное аналитич. пространство $(X, \widehat{G})$, обладающее следующими свойствами:



\begin{enumerate}

  \item любое комшактное аналитич. подмножество в $X$ конечно;



  \item любой компакт $K \subset X$ допускает такую открытую окрестность $W$ в $X$, что множество



\end{enumerate}



\[

\left\{x \in W \| f(x)\left|\leqslant \sup _{z \in K}\right| f(z) \quad \text { для всех } t \in G(X)\right\}

\]



компактно (с ла бая голомор фо а я вы пу кл $\mathbf{0}$ с т $\mathbf{~} \mathbf{)}$.



Комплексное многообразие $M$ тогда и только тогда является Ш. п., когда $M-$ III тейна многообразие. Комплексвое пространство является III. П. тогда и только тогда, когда этим свойством обладает его редукция. Всякое голоморфно выпуклое открытое подпространство в III. П. является III. П. Приведенное комшлексное пространство штейново тогда и только тогда, когда его нормализация есть II. П. Всякое замкнутое аналитич. подпространство в Ш. п., напр. в $\mathbb{C}^{n}$, есть ШI. п. Всякое конечномерное II. П. допускает собственное инъективное голоморфное отображение в нек-рое $\mathbb{C}^{n}$, регулярное в каждой неособой точке. Всякое неразветвленное накрытие III. П. есть ШІ. П. Прямое произведение двух ШІ. п. есть Ш. П. Тем же свойством обладает во многих случаях голоморфное расслоение, база и слой к-рого суть Ш. п. (напр., если структурной группоӥ является комплексная группа Ли с конечным числом связных компонент). Однако существуют примеры голоморфных расслоений со слоем $\mathbb{C}^{2}$ и базой $\mathbb{C}$, не являющихся многообразиями Штейна [2 .



\includegraphics[max width=\textwidth, center]{2023_07_09_417214854ba6de95322eg-1(1)}

( $X, 6)$. Тогда справедливы Картана теоремы:



A) Пространство $H^{0}(X, \mathscr{F})$ порождает слой $\mathscr{F}_{x}$ пучка \& в любой точке $x \in X$;



B) $H^{q}(X, \mathscr{J})=0$ для всех $q>0$.



Обратно, если $H^{1}(X, y)=0$ для любого когерентного пучка идеалов $\mathfrak{J} \subseteq \mathcal{G}$, то $X$-III. п. Область $D \subset \mathbb{C}^{n}$ является многообразием Штейна тогда и только



\begin{center}

\includegraphics[max width=\textwidth]{2023_07_09_417214854ba6de95322eg-1}

\end{center}



Из теорем Картана следует, что на III. П. всегда разрешима 1-я проблема Кузена, а если $H^{2}(X, \mathbb{Z})=0-$ то и 2-я проблема Кузена (см. Кузена проблемьі). На любом многообразии Штейна $X$ разрешима п ро бл е ма П у анк а р е, т. е. всякая мероморфнная функция представима в виде $\frac{f}{g}$, где $t, \quad g \in G(X), \quad g \neq 0$. Если при әтом $H^{2}(X, \mathbb{Z})=0$, то $f, g$ можно выбрать так, чтобы ростки $f_{x}$, $g_{x}$ в любой точке $x \in X$ были взаимно просты. Грушпа классов дивизоров неприводимого приведенного III. П. $X$ изоморфиа $H^{2}(X, \mathbb{Z})$. Для любого $n$-мерного III. п. $X$ группа гомологий $H_{q}(X, \mathbb{Z})=0$ для $q>n$, а группа $H_{n}(X, \mathbb{Z})$ не имеет кручения. Если $X$ - многообразие, то $X$ гомотопически эквивалентно $n$-мерному клеточному комплексу. С другой стороны, для любой счетной абелевой группы $G$ и для любого $q \geqslant 1$ существует область голоморфоности $D \subset \mathbb{C}^{2 q+3}$ такая, что $H_{q}(D, \mathbb{Z}) \cong$ $\cong G$.



Важное ваправление в теории II. п. связано с изучением плюрисубгармонич. функций на них (см. Леви проблема, Псевдовыпуклость $и$ псевдовогнутость). Основной результат здесь состоит в характеризации III. II. как пространства, на к-ром существует исчерпывающая его сильно 1-псевдовыпуклая функция.



Алгебры голоморфных функций $G(X)$ на III. п. $X$ (т. н. ш т е й н о в ы а л г е б р ы) обладают следующими свойствами. Для максимального идеала $I \subset G(X)$ эквивалентны условия: $I$ замкнут в $G(X)$ относительно топологии компактной сходимости; $I=\{f \in G(X) \mid f(x)=$ $=0\}$ для нек-рой точки $x \in X ; I$ конечно порожден. Если $X$ конечномерно, то каждый характер $\chi: \mathcal{O}(X) \longrightarrow \mathbb{C}$ имеет вид $\chi(f)=f(x)$ для нек-рой точки $x \in X$. Если $\left(X, \sigma_{X}\right),\left(Y, \widehat{\sigma}_{Y}\right)$ - два конечномерных III. І. с изоморфными алгебрами $\sigma_{X}(X) \cong \sigma_{Y}(Y)$, то $\left(X, \sigma_{X}\right) \cong\left(Y, \sigma_{Y}\right)$, причем любой изоморфизм $\hat{\sigma}_{X}(X) \rightarrow G_{Y}(Y)$ непрерывен и индуцируется нек-рым п1зоморфизмом $Y \rightarrow X$ комплексных пространств.



Большую роль в теории III. П. играет т. н. П р и нц и ІІ О к а, согласно к-рому многие задачи разрешимы на III. п. в классе аналитич. функций тогда и только тогда, когда они разрешимы в классе гладких непрерывных функций. Этому принципу удовлетворяет, напр., 2-я проблема Кузена. Более общим является следующее утверждение: классификация главных аналитических расслоений, базой к-рых является заданное приведенное III. п. $X$, а структурной группой заданная комплексная групша Ли $G$, совпадает с классификацией топологич. расслоений с той же базой и структурной группой. Совпадают также групшы связных компонент в групшах аналитических и непрерывных функций $X \rightarrow G$



Jum.: [1] G r a u e r t H., R e m m e r t R., Theorie der Steinschen Räume, B.- Hdlb - N. Y., 1977; [2] D e m a i 1[3] Итоги науки и техники. Алгебра. Топология. Геометрия, т. 11, M., 1974 , c. $125-51$; т. 15 , M., 1977 , c. $93-171$.



A. J. Онищик.



IIТЕЙНЕРА КРИВАЯ - плоская алгебраич. кривая 4-го порядка, к-рая описывается точкой окружности





\end{document}